\documentclass[11pt]{letter}
\usepackage[margin=1in]{geometry}
\usepackage[T1]{fontenc}
\usepackage{lmodern}
\usepackage{microtype}
\usepackage{enumitem}
\usepackage[hidelinks]{hyperref}

\setlist{itemsep=1pt,topsep=3pt}
\signature{Daniel Kushnir and Hao Wang}
\address{Daniel Kushnir (corresponding author) \quad \href{mailto:kushnir.daniel05@gmail.com}{kushnir.daniel05@gmail.com}\\
Hao Wang \quad \href{mailto:hoguewang@gmail.com}{hoguewang@gmail.com}\\
Rutgers University, New Brunswick, NJ, USA}
\date{October 30, 2026}

\begin{document}
\begin{letter}{Editors, Machine Learning Open Source Software Section\\
Journal of Machine Learning Research}

\opening{Dear Editors,}

We submit the manuscript \emph{TSFM-Lens: Controlled, Causal and Confirmatory Interpretability of Time-Series
Foundation Models} for consideration in the \textbf{machine learning open source software section} of the Journal of
Machine Learning Research.

Time-series foundation models such as TimesFM, the Chronos family and Sundial are widely used for zero-shot
forecasting, but the interpretability libraries built for language models do not transfer to them directly. Tokens
span between 1 and 96 time steps depending on the architecture, so positions cannot be compared across models; real
series carry no labels that could validate a probe; and the causal effects of internal features are small enough that
uncontrolled measurements produce confident false positives. TSFM-Lens is a Python toolkit built around these three
problems. Given a list of checkpoints and one configuration file, it measures a common time axis across architectures,
validates its probes on a leakage-audited synthetic benchmark with exact generative ground truth, tests sparse-autoencoder
concepts against matched random-direction nulls with reach controls, confirms registered claims once on a sealed held-out
split, and writes a single self-contained HTML report in which every claim states its evidence class. To our knowledge no
existing library provides these components for time-series models; the manuscript compares it with TransformerLens,
nnsight, pyvene, SAELens, Captum and recent sparse-autoencoder studies of individual forecasting models.

The information requested for submissions to this section is as follows.
\begin{enumerate}
\item \textbf{License.} The software is released under the MIT license. The license is stated in the \texttt{LICENSE}
file and in the package metadata, and every source file carries an SPDX license identifier. All dependencies (PyTorch,
NumPy, SciPy, scikit-learn, zarr, Plotly and others) are open source.
\item \textbf{Project web address.} \url{https://github.com/Wang-ML-Lab/TSFM-Lens}. Both packages are installable from
PyPI (\texttt{tsfm-lens} and \texttt{tsfm-benchmark}), and each release is archived on Zenodo
(\href{https://doi.org/10.5281/zenodo.23267462}{doi:10.5281/zenodo.23267462}).
\item \textbf{Version to be reviewed.} Version 1.1.0.
\item \textbf{User community.} The project is young and its community is small but active. The repository is public
under the Wang ML Lab organization at Rutgers, has 2 stars on GitHub, and several researchers use the toolkit in their
own work. The PyPI package \texttt{tsfm-lens} has been downloaded 128 times since its first release on October~9, 2026
(pypistats.org, excluding mirrors). The issue tracker is open, and the repository provides a contributing guide, issue
templates (including one for requesting a new model adapter) and a pull-request template.
\item \textbf{Quality and documentation.} The repository contains over 2{,}800 tests, which run in continuous
integration on Linux against pinned dependencies on every push and against the latest releases of the load-bearing
dependencies weekly; contributors are asked to show that each load-bearing assertion fails against a planted regression.
Documentation includes a quickstart (also as a Colab notebook), tutorial notebooks, a worked example, a guide and template
for adding a model, generated stage and glossary references, a changelog, and a complete seven-model example run with its
report and held-out confirmation.
\item \textbf{Prior publication.} The software has not been the subject of a peer-reviewed publication, and this
manuscript is not under review elsewhere. The scientific results obtained with the toolkit are reported in a separate
manuscript submitted to JMLR as a regular paper; that manuscript does not describe the software's design, and the
present submission does not report its scientific results.
\item \textbf{Consent of co-authors and conflicts of interest.} Both authors are aware of this submission and consent
to its review by JMLR. The authors have no competing interests, and the work received no external funding.
\end{enumerate}

\textbf{Use of AI tools.} The authors conceived the research questions, designed the software and its methodology,
directed its implementation, reviewed the code and checked every reported result. Claude (Anthropic) served as an AI
coding and writing assistant throughout: it generated much of the code and tests, ran experiments and helped draft and
edit the manuscript. The authors take full responsibility for the content.

Please direct correspondence to Daniel Kushnir at \href{mailto:kushnir.daniel05@gmail.com}{kushnir.daniel05@gmail.com}.
Thank you for considering our submission.

\closing{Sincerely,}

\end{letter}
\end{document}
